\documentclass{tufte-handout}

\title{Quartic Reciprocity Proof}

\author[Asher]{Asher Lucas-Cuddeback}

%\date{28 March 2010} % without \date command, current date is supplied

%\geometry{showframe} % display margins for debugging page layout

\usepackage{graphicx} % allow embedded images
  \setkeys{Gin}{width=\linewidth,totalheight=\textheight,keepaspectratio}
  \graphicspath{{graphics/}} % set of paths to search for images
\usepackage{amsmath}  % extended mathematics
\usepackage{amsthm}
\usepackage{booktabs} % book-quality tables
\usepackage{units}    % non-stacked fractions and better unit spacing
\usepackage{multicol} % multiple column layout facilities
\usepackage{lipsum}   % filler text
\usepackage{fancyvrb} % extended verbatim environments
  \fvset{fontsize=\normalsize}% default font size for fancy-verbatim environments

% Standardize command font styles and environments
\newcommand{\doccmd}[1]{\texttt{\textbackslash#1}}% command name -- adds backslash automatically
\newcommand{\docopt}[1]{\ensuremath{\langle}\textrm{\textit{#1}}\ensuremath{\rangle}}% optional command argument
\newcommand{\docarg}[1]{\textrm{\textit{#1}}}% (required) command argument
\newcommand{\docenv}[1]{\textsf{#1}}% environment name
\newcommand{\docpkg}[1]{\texttt{#1}}% package name
\newcommand{\doccls}[1]{\texttt{#1}}% document class name
\newcommand{\docclsopt}[1]{\texttt{#1}}% document class option name
\newenvironment{docspec}{\begin{quote}\noindent}{\end{quote}}% command specification environment
\newtheorem{thm}{Theorem}
\begin{document}

\maketitle% this prints the handout title, author, and date

Take \(f(x)=x^4-2\). If \(f\) splits modulo some \(p\) then the quotient of two of its roots
must be a square root of \(-1\) in \(\mathbb{Z}/p\mathbb{Z}\). 

\begin{thm}
  The splitting behavior of a polynomial \(f(x)\) modulo \(p\) is determined by congruence 
  conditions if and only if the roots are contained in a cyclotomic field.
\end{thm}

We can employ theorem 1 to get that \(p\equiv 1 \pmod{4}\). 


\begin{thm}[Fermat's Theorem on Square Primes]
  An odd prime, \(p\) can be expressed as \(p\equiv 1\pmod{4}\) if and only if
  \(p=a^2+b^2\).
\end{thm}

By Fermat's theorem, \(p=a^2+b^2\). Without loss of generality, assume that \(a\) is even and
\(b\) is odd. We find ourselves now in the ring of Gaussian integers. Here in \(\mathbb{Z}[i]\),
\(p\) is not prime. Instead, we can write \(p\) as \(p=\pi\overline{\pi}\) with \(\pi=a+bi\). 

\begin{thm}[2.2.3]
  \(x^4-2\) splits modulo \(p\) if and only if \(\pi\equiv 1,3,5,7 \pmod{8\mathbb{Z}[i]}\).
\end{thm}
This translates to \(b=8b_0\) for integer \(b_0\) so \(p=a^2+64b_0^2\).


\end{document}

